%% region CR Dynamics
% Draft scope: pool classified learner responses to characterize CR acquisition, timing, and extinction.
% todo ensure that learner classification and timing analyses use independent features or trials.









%!!! talk about this in Figure 4, when characterizing the CR in learners!
A comparison between the Delay and trace groups and the respective control fish revealed that only the groups that received the CS-US paired during training developed a consistent behavioural change in response to the CS during the training and testing phases 
During the first 10 training trials, larvae in the Delay condition showed a progressive reduction in movement vigor within the CS--US interval, emerging after several trials and strengthening with continued training. This conditioned response (CR) became increasingly temporally refined: it initially appeared a few seconds after CS onset, then gradually shifted earlier to align more closely with CS presentation. During early test trials, the suppression often extended beyond CS offset, but it gradually extinguished across testing, with vigor returning to baseline as the CS regained a neutral value.
Larvae trained with short temporal gaps--- fixed \SI{3}{\second} trace interval (3sTrace)---displayed robust associative learning comparable to that observed in the Delay condition. In both paradigms, the CR first appeared near CS offset, close to the expected US time. As training progressed, the CR shifted earlier to occur shortly after CS onset. During testing, the CR was evident in early trials but extinguished over time, leaving only a mild, startle-like response similar to controls.
When the trace interval was extended to \SI{10}{\second}, conditioning became substantially more difficult. Heatmaps of pooled scaled vigor revealed only a subtle suppression that developed after several training trials, localized primarily within the trace interval (between CS offset and US onset). This weak and delayed reduction---absent in control fish---suggests that some associative learning can occur even across this extended temporal gap, though with greatly diminished strength.






% Previous draft heading: CR Dynamics Reveal Precise Temporal Expectation of the US in Delay and Short-Trace Conditioning
\subsection{Pooled Learner Responses Characterize Conditioned-Response Dynamics}

To characterize how the CR developed and when it was expressed relative to the expected US, we combined two complementary analyses: (i) block-averaged scaled-vigor traces (10-trial blocks) and (ii) catch-trial traces (CS-only trials embedded in training).
% todo averaged...?

In the delay condition, CRs emerged rapidly and became both stronger and temporally precise with training. Block averages showed an initial suppression of vigor during the latter part of the CS in the earliest training block (train 1), which strengthened and shifted earlier over subsequent blocks. By late training, suppression began shortly after CS onset and peaked before the expected US time. Catch-trial traces confirmed this timing: CR suppression typically began within a few seconds of CS onset, persisted for about \SIrange{1}{2}{\second} after CS offset, and returned to baseline shortly after the expected US time (Fig. X). During early testing (CS alone), the CR remained robust and time-locked to the predicted US; across test blocks the response shortened and weakened, consistent with extinction (Fig. X).

The and 3sTrace groups displayed qualitatively similar dynamics to delay conditioning but with trace-specific structure. When trace intervals were short ( initially used a \SI{0.5}{\second} gap and was progressively extended; 3sTrace used a fixed \SI{3}{\second} gap), block averages showed a suppression that began during the CS and extended into the stimulus-free interval. With training, the suppression intensified and became focused so that pooled catch-trial data peaked within the trace interval at the predicted US time (Fig. X; Suppl. Fig. Y). These patterns indicate that larvae maintained a CS representation across the gap and developed accurate temporal expectation of US onset. As in delay conditioning, CRs in these groups persisted into early testing and then declined.

The 10sTrace group showed clearly weaker and more variable CRs. Block averages indicated only a modest suppression emerging late in training and persisting into the early test blocks; suppression was largely confined to the long trace interval, whereas scaled vigor during the CS itself stayed near baseline (Fig. X). Pooled catch-trial averages revealed a subtle, diffuse reduction in vigor across the extended trace window that often extended beyond the expected US time. Quantitatively, the 10sTrace group exhibited a statistically significant but attenuated reduction in normalized vigor during the early test phase compared with both pre-train baseline and controls (see Methods and Fig. X), consistent with associative learning that is imprecise in time when the CS--US interval is long.

Across conditions, two consistent themes emerged. First, catch trials isolated a temporally precise CR in delay and short-trace paradigms: suppression was anticipatory (appearing before the expected US) and returned to baseline shortly after the expected US time, demonstrating accurate temporal estimation. Second, single-fish responses were heterogeneous---some individuals showed strong, well-timed CRs while others did not---so pooled averages are complemented by exemplar single-trial traces (Suppl. Fig. Y). Finally, CRs diminished across test blocks in all learned groups, reflecting extinction when the CS was repeatedly presented without the US.











% endregion CR Dynamics
